\chapter{Inverse Problems in Biomechanics}
\label{ch:inverse-problems}

\epigraph{%
  In biomechanics, we almost never measure forces directly.
  We measure motion and infer the forces that caused it.
}{}

\section{The Inverse Dynamics Pipeline}

Inverse dynamics is the most widely used computational tool in biomechanics: given
measured kinematics $\bm{q}(t)$, $\dot{\bm{q}}(t)$, $\ddot{\bm{q}}(t)$ and external
forces $\bm{F}_{\text{ext}}(t)$, compute the net joint torques. The complete pipeline,
illustrated in Figure~\ref{fig:ch6:id-pipeline}, transforms raw experimental data into
mechanistic estimates of the forces acting on the system:
\begin{equation}
  \bm{\tau}(t) = M(\bm{q})\ddot{\bm{q}} + C(\bm{q}, \dot{\bm{q}})\dot{\bm{q}} + g(\bm{q})
  - J_c^T(\bm{q}) \bm{F}_{\text{ext}}.
  \label{eq:ch6:inverse-dynamics}
\end{equation}

\begin{center}
\begin{tikzpicture}[scale=0.9, >=Stealth, node distance=2.2cm]
  \node[draw, thick, minimum width=2.5cm, fill=blue!10] (collect) {Data Collection\\(markers, forces)};
  \node[draw, thick, minimum width=2.5cm, fill=blue!10, right of=collect] (process) {Marker\\Processing};
  \node[draw, thick, minimum width=2.5cm, fill=green!10, right of=process] (ik) {Inverse\\Kinematics};
  \node[draw, thick, minimum width=2.5cm, fill=green!10, right of=ik] (id) {Inverse\\Dynamics};
  \node[draw, thick, minimum width=2.5cm, fill=red!10, right of=id] (muscle) {Muscle Force\\Estimation};

  \draw[->, thick] (collect) -- (process);
  \draw[->, thick] (process) -- (ik);
  \draw[->, thick] (ik) -- (id);
  \draw[->, thick] (id) -- (muscle);

  \node[below, font=\small] at (collect.south) {Observables};
  \node[below, font=\small] at (ik.south) {Kinematics};
  \node[below, font=\small] at (id.south) {Net torques};
  \node[below, font=\small] at (muscle.south) {Individual forces};
\end{tikzpicture}
\end{center}

The standard pipeline comprises five stages:
\begin{enumerate}
  \item \textbf{Data Collection}: Motion capture markers + force plates
  \item \textbf{Marker Processing}: Filtering, gap-filling, labeling
  \item \textbf{Inverse Kinematics}: Fit skeleton model to marker data
  \item \textbf{Inverse Dynamics}: Compute joint torques from RNEA
  \item \textbf{Muscle Force Estimation}: Distribute torques among muscles
\end{enumerate}

\subsection{Inverse Kinematics}

Given marker positions $\bm{m}_i^{\text{meas}}(t) \in \Reals^3$ for $i = 1, \ldots, N_m$
markers, inverse kinematics solves:
\begin{equation}
  \bm{q}^*(t) = \argmin_{\bm{q}} \sum_{i=1}^{N_m} w_i
  \norm{\bm{m}_i^{\text{model}}(\bm{q}) - \bm{m}_i^{\text{meas}}(t)}^2,
  \label{eq:ch6:ik-optimization}
\end{equation}
where $\bm{m}_i^{\text{model}}(\bm{q})$ is the predicted marker position from the
skeleton model and $w_i$ are weighting factors.

\begin{lstlisting}[language=Python, caption={Simplified inverse kinematics solver.}]
import numpy as np
from scipy.optimize import minimize

def inverse_kinematics_frame(
    marker_positions_measured: np.ndarray,
    marker_positions_model_func,
    q_initial: np.ndarray,
    weights: np.ndarray | None = None
) -> np.ndarray:
    """Solve inverse kinematics for a single frame.

    Parameters
    ----------
    marker_positions_measured : np.ndarray, shape (n_markers, 3)
        Measured marker positions.
    marker_positions_model_func : callable
        Function q -> (n_markers, 3) predicted marker positions.
    q_initial : np.ndarray, shape (n_dof,)
        Initial guess for joint angles.
    weights : np.ndarray, shape (n_markers,), optional
        Per-marker weights.

    Returns
    -------
    np.ndarray, shape (n_dof,)
        Optimal joint angles.
    """
    n_markers = marker_positions_measured.shape[0]
    if weights is None:
        weights = np.ones(n_markers)

    def objective(q: np.ndarray) -> float:
        m_model = marker_positions_model_func(q)
        residuals = m_model - marker_positions_measured
        return float(np.sum(weights[:, None] * residuals**2))

    result = minimize(objective, q_initial, method='L-BFGS-B')
    return result.x
\end{lstlisting}

\subsection{Bottom-Up vs.\ Top-Down Inverse Dynamics}

\textbf{Bottom-up}: start from the ground contact (where forces are measured by
force plates) and propagate upward through the kinematic chain. Most accurate when
force plate data is available.

\textbf{Top-down}: start from the most distal segment and work inward. Useful when
ground reaction forces are not available (e.g., swimming, throwing).

\section{The Muscle Redundancy Problem}

Inverse dynamics gives net joint torques $\bm{\tau} \in \Reals^n$, but we want
individual muscle forces $\bm{F} \in \Reals^p$ with $p > n$ (typically $p = 3n$
to $5n$). This is an underdetermined system:
\begin{equation}
  R(\bm{q}) \bm{F} = \bm{\tau}, \quad \bm{F} \geq 0.
  \label{eq:ch6:redundancy}
\end{equation}

\subsection{Static Optimization}

The most common approach minimizes a cost function:
\begin{equation}
  \min_{\bm{F}} \sum_{i=1}^{p} \left(\frac{F_i}{F_{0,i}}\right)^\alpha
  \quad \text{subject to} \quad R(\bm{q})\bm{F} = \bm{\tau},\;
  0 \leq F_i \leq a_i \cdot f_L(\tilde{\ell}_i) \cdot F_{0,i},
  \label{eq:ch6:static-opt}
\end{equation}
where $\alpha = 2$ (minimize sum of squared activations) or $\alpha = 3$ (minimize
metabolic energy proxy). The constraint $F_i \leq a_i \cdot f_L \cdot F_{0,i}$ enforces
the force-length and activation constraints from Chapter~3.

\subsection{Computed Muscle Control (CMC)}

CMC \cite{Thelen2003} combines forward dynamics with feedback control. It uses a
PD controller to track the measured kinematics, with the controller output being
the desired muscle excitations:
\begin{equation}
  \ddot{\bm{q}}_{\text{des}} = \ddot{\bm{q}}_{\text{meas}} + K_v(\dot{\bm{q}}_{\text{meas}} - \dot{\bm{q}})
  + K_p(\bm{q}_{\text{meas}} - \bm{q}),
  \label{eq:ch6:cmc}
\end{equation}
where $K_v$ and $K_p$ are velocity and position feedback gains.

\section{Parameter Estimation}

Many biomechanical parameters (segment masses, muscle strengths, joint axis locations)
are not directly measurable and must be estimated from motion data.

\subsection{Bayesian Framework}

The parameter estimation problem is naturally Bayesian:
\begin{equation}
  p(\bm{\theta} | \bm{y}) \propto p(\bm{y} | \bm{\theta}) \cdot p(\bm{\theta}),
  \label{eq:ch6:bayes}
\end{equation}
where $\bm{\theta}$ are model parameters, $\bm{y}$ are observations, $p(\bm{y} | \bm{\theta})$
is the likelihood, and $p(\bm{\theta})$ is the prior (from literature values and
physical constraints).

\section*{Exercises}

\begin{enumerate}
  \item \textbf{Inverse Dynamics.}
        For a 2-segment pendulum model of the arm, compute the shoulder and elbow
        torques required during a reaching movement given $\bm{q}(t) = (\sin(\pi t/T), \pi t / (2T))$.

  \item \textbf{Muscle Redundancy.}
        For a single-DOF elbow with 3 muscles (biceps, brachialis, brachioradialis),
        solve the static optimization for a given elbow torque. Show that the solution
        depends on the cost function exponent~$\alpha$.

  \item \textbf{(Programming.)} Implement the static optimization for a 2-DOF,
        6-muscle system. Compare $\alpha = 2$ and $\alpha = 3$ solutions.

  \item \textbf{(Programming.)} Implement an inverse kinematics solver for a
        3-segment arm model with 7 markers. Test with synthetic marker data.
\end{enumerate}
